% Vietnamese translation for OI Public Library
% Source-File: IOI中国国家候选队论文/2005/杨思雨/试题/取石子游戏.doc
\documentclass[11pt]{article}
\usepackage{oipl-vn}

\title{Trò chơi lấy đá}
\author{SHTSC 2002}
\date{}

\begin{document}
\maketitle

\origtitle{STONE}
\sourcepath{IOI Candidate Team Papers/2005/Yang Siyu/problems/STONE.doc}

\section{Phát biểu}

Có hai đống đá, số lượng đá trong hai đống là tùy ý và có thể khác nhau. Hai
người chơi luân phiên lấy đá. Mỗi lượt có hai kiểu lấy:
\begin{itemize}
  \item chọn một trong hai đống và lấy đi một số đá dương bất kỳ;
  \item lấy đồng thời cùng một số đá dương từ cả hai đống.
\end{itemize}
Người lấy hết toàn bộ số đá cuối cùng là người thắng.

Cho số đá ban đầu của hai đống. Nếu đến lượt bạn đi trước và cả hai người đều
chơi tối ưu, hãy xác định cuối cùng bạn thắng hay thua.

\section{Dữ liệu vào và ra}

Tệp vào gồm nhiều dòng, mỗi dòng mô tả một trạng thái ban đầu. Mỗi dòng chứa
hai số nguyên không âm \(a\) và \(b\), là số đá trong hai đống, với
\[
  a,b\le 1\,000\,000\,000.
\]

Tệp ra cũng gồm số dòng tương ứng. Với mỗi trạng thái, in ra \(1\) nếu người
đi trước thắng, ngược lại in ra \(0\).

\section{Ví dụ}

\paragraph{Ví dụ 1.}
\begin{verbatim}
Input
6 10

Output
0
\end{verbatim}

\paragraph{Ví dụ 2.}
\begin{verbatim}
Input
2 1
8 4
4 7

Output
0
1
0
\end{verbatim}

\end{document}
